\section{Clasificación de los métodos de Meta-RL}

\subsection{Planteamiento típico de una tarea de Meta-RL}

El ejemplo clásico de Meta-RL es maze navigation: primero se recopila una pequeña cantidad de experiencia exploratoria en un laberinto nuevo y luego esa experiencia se usa para resolver ese laberinto. La dificultad propia de este caso es que el exploration-exploitation trade-off cambia: las primeras interacciones no solo deben obtener reward en el momento, sino también identificar de qué tarea se trata.

\begin{figure}[H]
\centering
\includegraphics[width=0.9\textwidth]{slides-images/slide-10.jpg}
\caption{Ejemplo de Meta-RL: primero se explora el laberinto nuevo y después se usa poca experiencia para resolver la tarea}
\end{figure}

\begin{figure}[H]
\centering
\includegraphics[width=0.9\textwidth]{slides-images/slide-11.jpg}
\caption{Diferencia entre meta-train y meta-test: durante el entrenamiento se aprende a adaptarse con eficiencia; durante la prueba se enfrenta una tarea nueva y se domina rápidamente}
\end{figure}

\subsection{Formulación más formal del problema}

Esta clase presentó un formal setup muy importante: Meta-RL no es un solo MDP, sino una \textbf{distribución de tareas} $p(\mathcal{T})$. En la fase de entrenamiento se extraen muchas tareas de ella y se aprende un mecanismo de adaptación; en la fase de prueba, ante una tarea nueva de la misma distribución, solo se dispone de poca experiencia.

\begin{figure}[H]
\centering
\includegraphics[width=0.9\textwidth]{slides-images/slide-13.jpg}
\caption{Definiciones más formales de multi-task, transfer y meta-learning}
\end{figure}

\begin{importantbox}{Entrada y salida de Meta-RL}
Tomando como ejemplo la versión episódica:
\begin{itemize}
    \item Entrada: una pequeña cantidad de rollout / experience $D_{\text{train}}$ proveniente de cierta tarea
    \item Salida: una política ajustada según $D_{\text{train}}$, $\pi(a\mid s, D_{\text{train}})$
\end{itemize}
Es decir, el task ID que en el RL multitarea tradicional se proporciona de forma explícita, en Meta-RL se sustituye por un breve tramo de experiencia de interacción.
\end{importantbox}

\begin{figure}[H]
\centering
\includegraphics[width=0.9\textwidth]{slides-images/slide-14.jpg}
\caption{Entrada y salida de Meta-RL: el contexto de experiencia sustituye al task descriptor explícito}
\end{figure}

\subsection{Método 1: métodos basados en contexto (Context-Based)}

La idea central es introducir la experiencia de la tarea nueva como \textbf{contexto} en la red de política, para que la red infiera la tarea a partir del contexto y tome decisiones.

\begin{importantbox}{Política con contexto}
$$\pi_\theta(a | s, c)$$
donde $c = \{(s_1, a_1, r_1, s_1'), \ldots, (s_k, a_k, r_k, s_k')\}$ es la poca experiencia recopilada en la tarea nueva. La red de política debe inferir las características de la tarea a partir de esa experiencia.

Arquitecturas comunes:
\begin{itemize}
    \item Codificar el contexto como estado oculto mediante una RNN/LSTM
    \item Codificar el contexto mediante el mecanismo de atención de un Transformer
    \item Codificar el contexto mediante un codificador de conjuntos (permutation-invariant)
\end{itemize}
\end{importantbox}

\begin{knowledgebox}{RL\textsuperscript{2} (Learning to Reinforcement Learn)}
RL$^2$ es uno de los primeros métodos de context-based meta-RL. Trata todo el proceso de RL como un POMDP más grande, en el que el estado oculto de la RNN codifica la «creencia» sobre la tarea actual. El algoritmo de RL externo entrena la RNN para que, en el nivel interno, se adapte rápidamente a tareas nuevas.
\end{knowledgebox}

\subsection{Ciclo de entrenamiento de Black-box Meta-RL}

Este curso se centra en black-box meta-RL: no se escribe de forma explícita una fórmula de actualización bayesiana, sino que se usa directamente una red neuronal con memoria para que el «cómo adaptarse a partir de la experiencia» quede aprendido en los parámetros.

\begin{figure}[H]
\centering
\includegraphics[width=0.9\textwidth]{slides-images/slide-17.jpg}
\caption{Algoritmo básico de Black-box Meta-RL: se recopilan varias episode por tarea y luego se actualiza de manera conjunta la política con memoria}
\end{figure}

\begin{knowledgebox}{En qué se distingue de una recurrent policy común}
A primera vista, black-box meta-RL se parece mucho a «RL con una RNN», pero el curso destacó dos condiciones adicionales:
\begin{enumerate}
    \item el reward normalmente entra a la red como parte de la entrada;
    \item el estado oculto debe conservarse a lo largo de varias episode de una misma tarea.
\end{enumerate}
Por eso, lo que la red recuerda no es solo una trayectoria individual, sino «lo que ya he observado en esta tarea».
\end{knowledgebox}

\subsection{Arquitecturas y optimizadores comunes}

El espacio de implementación de black-box meta-RL es amplio: puede usarse RNN, attention o Transformer, y también combinarse con distintos optimizadores de RL del outer-loop, como TRPO/A3C o SAC.

\begin{figure}[H]
\centering
\includegraphics[width=0.9\textwidth]{slides-images/slide-19.jpg}
\caption{Elección de arquitectura y optimizador en Black-box Meta-RL: RNN, attention, off-policy context variable, etc.}
\end{figure}

\subsection{Método 2: métodos basados en gradiente (Gradient-Based)}

\begin{importantbox}{MAML (Model-Agnostic Meta-Learning)}
La idea central de MAML es aprender un buen conjunto de parámetros iniciales $\theta$, de modo que, para cualquier tarea nueva, basten unos pocos pasos de gradiente para adaptarse.

\textbf{Actualización interna} (adaptación a la tarea):
$$\theta_i' = \theta - \alpha \nabla_\theta \mathcal{L}_{\mathcal{T}_i}(\theta)$$

\textbf{Actualización externa} (metaaprendizaje):
$$\theta \leftarrow \theta - \beta \nabla_\theta \sum_i \mathcal{L}_{\mathcal{T}_i}(\theta_i')$$

Punto clave: el objetivo de la optimización externa es que los parámetros \textbf{ya adaptados} $\theta_i'$ tengan buen desempeño en la tarea $\mathcal{T}_i$, no los parámetros originales $\theta$.
\end{importantbox}

\begin{warningbox}{Costo computacional de MAML}
MAML necesita calcular el gradiente de un «gradiente» (derivadas de segundo orden), lo que implica un costo computacional considerable. Aunque es posible reducirlo con aproximaciones de primer orden (como FOMAML o Reptile), esto sacrifica cierta precisión. En RL, como el gradiente de política ya tiene una varianza alta por sí mismo, la estabilidad del entrenamiento de MAML es un reto.
\end{warningbox}

\subsection{Comparación de los dos tipos de métodos}

\begin{table}[H]
\centering
\begin{tabular}{lcc}
\toprule
& \textbf{Context-Based} & \textbf{Gradient-Based} \\
\midrule
Forma de adaptación & Propagación hacia adelante (inferencia) & Descenso de gradiente \\
Velocidad de adaptación & Muy rápida (una sola pasada hacia adelante) & Requiere varios pasos de gradiente \\
Capacidad expresiva & Limitada por la capacidad de la red & En teoría permite una adaptación exacta \\
Dificultad de entrenamiento & Relativamente fácil & Requiere optimización de segundo orden \\
Métodos típicos & RL$^2$, PEARL & MAML, ProMP \\
\bottomrule
\end{tabular}
\caption{Comparación de los dos tipos de métodos de Meta-RL}
\end{table}

\subsection{Casos de Black-box Meta-RL}

El curso presentó dos casos muy representativos. El primero es la navegación visual en laberintos: el modelo hace meta-train en una gran cantidad de laberintos pequeños y, durante la prueba, debe explorar rápidamente laberintos que nunca ha visto y encontrar la salida. El segundo caso está más cerca del contexto actual de los modelos grandes: se considera cada problema de matemáticas como una tarea distinta y se optimiza el test-time compute, para que el modelo, con un presupuesto fijo de tokens, intente, verifique y elija estrategias con mayor eficiencia.

\begin{figure}[H]
\centering
\includegraphics[width=0.9\textwidth]{slides-images/slide-20.jpg}
\caption{Caso 1 de Meta-RL: navegación visual en laberintos, para comprobar si el modelo forma rápidamente una política en un laberinto nuevo}
\end{figure}

\begin{figure}[H]
\centering
\includegraphics[width=0.9\textwidth]{slides-images/slide-22.jpg}
\caption{Caso 2 de Meta-RL: considerar distintos problemas de matemáticas como tareas y optimizar el test-time compute}
\end{figure}

\subsection{Relación entre Meta-RL y las políticas multitarea}

El ponente dedicó una de las slides a explicar un punto que se presta a confusión: la política multitarea depende de un task ID explícito $z_i$, mientras que Meta-RL usa \textbf{la propia experiencia} como task identifier. En otras palabras, la política multitarea equivale a «yo te digo qué tipo de tarea es ahora», y Meta-RL a «infiérelo tú mismo a partir de la interacción».

\begin{figure}[H]
\centering
\includegraphics[width=0.9\textwidth]{slides-images/slide-24.jpg}
\caption{Cuando el task ID no está en el estado, la política debe inferir a partir de la experiencia en qué MDP se encuentra}
\end{figure}

\subsection{Resumen de la sección}

Meta-RL tiene dos grandes líneas técnicas: los métodos context-based logran una adaptación implícita al codificar la experiencia, y los métodos gradient-based logran un ajuste fino rápido al aprender una buena inicialización.
